\documentclass[11pt,a4paper]{article}
\usepackage{curso}
\begin{document}
\portada{Lección}{1.5}{El principio de Landauer}

\noindent\textbf{Unidad 1 --- Termodinámica e información.}\quad
\textbf{Objetivos:} (i) explicar por qué borrar un bit tiene un costo mínimo $\kB T\ln 2$;
(ii) derivar la cota con la segunda ley y con la energía libre; (iii) calcular su valor y compararlo
con dispositivos reales; (iv) distinguir reversibilidad lógica y termodinámica.
\textbf{Prerrequisitos:} lecciones 1.1 (energía libre), 1.3 (entropía de Shannon/Gibbs) y 1.4 (Szilard).

\section{¿Por qué importa?}
Toda la energía que consume un procesador termina como calor. La pregunta de Landauer (1961)
\cite{landauer} es si existe un costo \emph{mínimo}, impuesto por las leyes de la física y no por la
ingeniería, para procesar información. La respuesta es sutil: la mayor parte de las operaciones
lógicas \emph{podrían} hacerse sin disipación mínima, pero hay una que no: \textbf{borrar}.

\begin{intuicion}
Una memoria de un bit es un sistema físico con dos regiones distinguibles del espacio de estados
(«0» y «1»). Borrar significa llevarla a «0» \emph{sin importar} dónde estaba. Eso es fusionar dos
regiones en una: el volumen de espacio de fases accesible se reduce a la mitad. Comprimir un
sistema en contacto con un baño térmico cuesta trabajo, y ese trabajo se va al baño como calor.
\end{intuicion}

\section{El modelo de una partícula}
Considere un gas de una sola partícula en una caja de volumen $V$ en contacto con un baño a
temperatura $T$. Una pared divide la caja: izquierda $=0$, derecha $=1$. Para borrar, se retira la
pared y se comprime isotérmica y cuasiestáticamente con un pistón hasta $V/2$. Con la ecuación de
estado del gas ideal de una partícula, $pV=\kB T$, el trabajo realizado es
\begin{equation}
W = -\int_V^{V/2} p\,\dd V' = \kB T \int_{V/2}^{V}\frac{\dd V'}{V'} = \kB T\ln 2 .
\end{equation}
Como la energía interna de un gas ideal a $T$ constante no cambia, $Q=W=\kB T\ln 2$ sale al baño.

\section{Derivación general}
\subsection{Ruta 1: entropía y segunda ley}
\begin{definicion}[ (entropía de Gibbs)]
Para una distribución $p_i$ sobre los estados lógicos/físicos de la memoria,
$S=-\kB\sum_i p_i\ln p_i$. Con $H$ la entropía de Shannon en bits, $S=\kB\ln 2\,H$.
\end{definicion}
Antes del borrado (bit desconocido) $p=(\tfrac12,\tfrac12)$ y $S_{\text{antes}}=\kB\ln2$. Después
$p=(1,0)$ y $S_{\text{después}}=0$. La memoria pierde $\Delta S_{\text{mem}}=-\kB\ln2$. La segunda ley
aplicada a memoria $+$ entorno exige $\Delta S_{\text{mem}}+\Delta S_{\text{ent}}\ge0$. Un baño a
temperatura $T$ que recibe un calor $Q$ aumenta su entropía en $Q/T$. Por lo tanto
\begin{resultado}[ (cota de Landauer)]
\begin{equation}
Q_{\text{dis}}\;\ge\;\kB T\ln 2\qquad\text{por bit borrado (bit inicialmente uniforme).}
\end{equation}
Para un bit con distribución inicial $p$ que se reinicia a un estado conocido:
$Q_{\text{dis}}\ge \kB T\ln 2\; h(p)$, con $h(p)=-p\log_2p-(1-p)\log_2(1-p)$.
\end{resultado}

\subsection{Ruta 2: energía libre}
Para un proceso isotérmico, la segunda ley se escribe $W\ge\Delta F$ con $F=U-TS$. En un pozo doble
simétrico (los estados 0 y 1 tienen la misma energía) se cumple $\Delta U=0$, de modo que
$\Delta F=-T\Delta S=\kB T\ln 2$. Por la primera ley, $Q = W-\Delta U = W\ge \kB T\ln2$. Ambas rutas
coinciden. (La segunda ruta muestra además que, si los estados tuvieran energías distintas, parte del
costo podría compensarse con $\Delta U$; lo que es universal es la parte entrópica.)

\section{El borrado en un pozo doble (simulación)}
Una memoria realista es una partícula sobreamortiguada en un potencial
$V(x,t)=E_b\,[x^4-2s(t)x^2]+c(t)\,x$ (forma $ax^4-bx^2+cx$ del laboratorio L2) que obedece la dinámica de Langevin
$\dot x=-\mu\,\partial_xV+\sqrt{2\mu\kB T}\,\xi(t)$. Con $s=1$, $c=0$ hay dos pozos en $x=\pm1$ separados
por una barrera $E_b$. El protocolo de borrado es:
\begin{enumerate}[nosep]
\item bajar la barrera ($s:1\to0$): los dos estados se fusionan;
\item inclinar ($c:0\to c_{\max}$): toda la probabilidad se desplaza a $x<0$;
\item subir la barrera ($s:0\to1$) y quitar la inclinación ($c\to0$).
\end{enumerate}
El trabajo de una trayectoria se calcula con la definición de Sekimoto,
$W=\int_0^\tau \partial_tV(x_t,t)\,\dd t$ (Unidad 3.1). En las animaciones de esta lección usamos
$E_b=c_{\max}=8\,\kB T$, $\mu=1$, paso $\Delta t=10^{-3}$ y 4000 trayectorias por punto. Resultados
\evS:
\begin{center}
\begin{tabular}{lccccccc}\toprule
$\tau$ & 1 & 2 & 5 & 10 & 20 & 50 & 100\\\midrule
$\langle W\rangle/\kB T$ & 4,30 & 2,61 & 1,49 & 1,12 & 0,90 & 0,78 & 0,73\\\bottomrule
\end{tabular}
\end{center}
El trabajo medio se acerca a $\ln 2\approx0{,}693$ \emph{por arriba}, con un exceso que decae
aproximadamente como $a/\tau$ ($a\approx4\,\kB T$ en estas unidades). Para $\tau$ muy grande aparece un
efecto secundario: la barrera finita permite saltos térmicos de regreso (fallo de borrado), y la
cota relevante pasa a ser $\kB T\ln2\,[1-h(p_{\text{fallo}})]$.

\section{Ejemplo numérico}
\begin{ejemplo}
Con $\kB=1{,}380649\times10^{-23}$ J/K (valor exacto del SI) y $T=300$ K:
$\kB T\ln2=1{,}380649\times10^{-23}\cdot300\cdot0{,}693147\approx\SI{2.871e-21}{J}\approx\SI{0.0179}{eV}$.
A $T=4$ K la cota es 75 veces menor, $\approx\SI{3.83e-23}{J}$.
El pensum cita para un chip tipo Z1 una energía por actualización de pbit $e_{\text{update}}\le
\SI{7.4e-14}{J}$ \evP; el cociente es $7{,}4\times10^{-14}/2{,}87\times10^{-21}\approx2{,}6\times10^{7}$.
\end{ejemplo}

\section{Reversibilidad lógica y termodinámica}
Una función lógica es \emph{reversible} si es inyectiva (de la salida se reconstruye la entrada).
NOT es una biyección; AND envía $00,01,10\mapsto0$. Con entradas uniformes, la entropía de entrada es
2 bits y la de salida es $h(1/4)\approx0{,}811$ bits: se pierden $\approx1{,}19$ bits, que cuestan al
menos $1{,}19\,\kB T\ln2$ si esa información se descarta en el entorno.
Bennett \cite{bennett} mostró que cualquier cómputo puede hacerse de forma lógicamente reversible:
\emph{calcular} guardando la historia, \emph{copiar} el resultado y \emph{descalcular} para devolver los
registros auxiliares a su estado inicial. El borrado inevitable se reduce a lo que se quiera
desechar de verdad. La reversibilidad \emph{termodinámica} (disipación $\to0$) exige además
operar cuasiestáticamente, es decir, lentamente.

\section{Aplicaciones y contexto}
\begin{aplicacion}
\begin{itemize}[nosep]
\item \textbf{Confiabilidad.} La tasa de error térmico de un bit con barrera $E_b$ escala como
$e^{-E_b/\kB T}$; para errores de $10^{-17}$ se necesita $E_b\approx 39\,\kB T$. Los márgenes de ruido de
la lógica digital ya sitúan cada operación muy por encima de $\kB T\ln2$.
\item \textbf{Comunicación.} Cargar un cable de capacidad $C$ a voltaje $V$ disipa del orden de $CV^2$,
independiente de Landauer (Unidad 0.1).
\item \textbf{Cómputo termodinámico.} El objetivo no es alcanzar Landauer sino dejar de pagar por
suprimir el ruido: usarlo como fuente de aleatoriedad para muestrear (Unidades 4--5).
\item \textbf{Laboratorio L2.} Reproducir la tabla anterior, el histograma de $W$ y la curva
$\langle W\rangle(\tau)$; reto: cerrar el ciclo información--trabajo con la máquina de Szilard.
\end{itemize}
\end{aplicacion}

\section{Conexiones}
El principio resuelve el demonio de Maxwell (lección 1.4): el demonio debe borrar su memoria y ese
borrado paga lo ganado. En la Unidad 3 se formaliza el trabajo y el calor por trayectoria, se
demuestra la igualdad de Jarzynski y se cuantifica el costo de borrar rápido (límites de velocidad,
longitud termodinámica). La Unidad 3.7 trata el reinicio y la lectura sin errores, y la Unidad 5
compara energías reales de pbits con esta cota.

\begin{error}
\begin{itemize}[nosep]
\item \emph{«Toda operación cuesta $\kB T\ln2$».} Solo el borrado lógicamente irreversible tiene costo mínimo.
\item \emph{«El costo se paga al medir».} En la formulación de Landauer--Bennett se paga al borrar la memoria.
\item \emph{«Borrar cualquier bit cuesta $\kB T\ln2$».} El costo mínimo es $\kB T\ln2\cdot h(p)$; un bit conocido se reinicia gratis.
\item \emph{«Landauer limita a los chips actuales».} Hoy dominan los márgenes de ruido y el movimiento de datos (factor $\sim10^{7}$).
\item \emph{«La cota se alcanza con cualquier protocolo».} Solo en el límite cuasiestático; a $\tau$ finito hay exceso $\propto1/\tau$.
\end{itemize}
\end{error}

\section{Resumen}
\begin{itemize}[nosep]
\item Borrar un bit uniforme reduce la entropía de la memoria en $\kB\ln2$.
\item La segunda ley obliga a disipar $Q\ge\kB T\ln2\approx\SI{2.87e-21}{J}$ a 300 K.
\item Solo pagan las operaciones lógicamente irreversibles; Bennett muestra cómo evitarlas.
\item El límite solo se alcanza con protocolos lentos.
\end{itemize}

\section{Preguntas de verificación}
\begin{enumerate}[nosep]
\item ¿Qué cambia en la cota si la memoria está a 77 K (nitrógeno líquido)?
\item ¿Por qué el borrado de un bit conocido no necesita disipar calor?
\item En el protocolo del pozo doble, ¿en qué paso se pierde la información?
\end{enumerate}

\begin{thebibliography}{9}\small
\bibitem[8]{landauer} R. Landauer, Irreversibility and heat generation in the computing process, \emph{IBM J. Res. Dev.} 5, 183 (1961).
\bibitem[9]{bennett} C. H. Bennett, Notes on Landauer's principle, reversible computation and Maxwell's demon, 2003. \url{https://arxiv.org/abs/physics/0210005}
\bibitem[10]{wolpert} D. H. Wolpert, The stochastic thermodynamics of computation, \emph{J. Phys. A} 52, 193001 (2019).
\bibitem[2]{peliti} L. Peliti y S. Pigolotti, \emph{Stochastic Thermodynamics: An Introduction}, Princeton University Press, 2021.
\end{thebibliography}
\noindent{\footnotesize Numeración [n] según la bibliografía del pensum (\texttt{references/bibliografia.md}).}
\end{document}
